\section{Thesis Contributions} \label{sec:14-introduction-contributions}

\begin{foldable}
  The three restricted settings identified in \S\ref{sec:12-introduction-gap} and examined in \S\ref{sec:13-introduction-fidelity} — few-shot \ac{KD}, data-free \ac{KD}, and text dataset distillation — each impose a distinct form of information scarcity on the distillation process.
  This thesis addresses all three through a common design philosophy: treat distributional fidelity as a first-class design objective, derive practical constraints from the deployment setting, and apply principled quality control to any synthetic or selected data that enters the training pipeline.
  The three contributions are as follows.
\end{foldable}

\begin{enumerate}

  \item \textbf{DivBFKD — Diversity-Aware Black-Box Few-Shot Knowledge Distillation (Chapter~\ref{ch:30-research01}).}
    \begin{foldable}
      DivBFKD addresses the setting in which only a small number of real training images are available and the teacher is accessible only as a black box, returning final predictions rather than internal representations.
      The method trains a \ac{WGAN} to synthesise additional training images, augmenting this process with a high-confidence image scheme that identifies and retains synthetic samples whose teacher-assigned confidence exceeds class-specific adaptive thresholds computed via the $q$-quantile of the score distribution.
      This threshold mechanism acts as a quality gate, ensuring that only sufficiently informative and class-coherent synthetic images enter the distillation pipeline, directly counteracting the semantic narrowness that arises when few real examples are available.
      Evaluated across seven image classification benchmarks — MNIST, Fashion-MNIST, SVHN, CIFAR-10, CIFAR-100, Tiny-ImageNet, and Imagenette — DivBFKD establishes state-of-the-art performance on all seven, achieving a gain of over 17\% relative to standard \ac{KD} on CIFAR-10 and closing to within 2\% of a fully supervised student upper bound on simpler datasets.
      This work was published at ECML PKDD 2024.
    \end{foldable}

  \item \textbf{DIPKD — Diverse Image Priors for Black-Box Data-Free Knowledge Distillation (Chapter~\ref{ch:40-research02}).}
    \begin{foldable}
      DIPKD addresses the strictest image-domain setting: the student has access to neither real training data nor teacher internals, receiving only top-1 hard labels in response to queries.
      The method is organised as a three-phase pipeline.
      In the first phase, a diverse pool of synthetic images is constructed by applying hierarchical noise injection, nonlinear photometric transforms, and semantic CutMixing, producing a set of candidate images that covers a broad region of the input space.
      In the second phase, a lightweight primer student $S_0$ is trained on these candidates under a contrastive objective that explicitly maximises inter-class separation, sharpening the diversity of the pool before full distillation begins.
      In the third phase, the trained primer student provides soft distillation signals — approximate logits — to supplement the hard labels returned by the teacher, recovering much of the information lost by black-box access.
      DIPKD achieves state-of-the-art results on 11 of 12 benchmarks, with a gain of over 21\% on Imagenette relative to the best prior baseline, and demonstrates strong performance on medical imaging datasets from MedMNIST.
      This work has been submitted to ICCCI 2026.
    \end{foldable}

  \item \textbf{TAKE — Trajectory-Aware Knowledge Estimation for Text Dataset Distillation (Chapter~\ref{ch:50-research03}).}
    \begin{foldable}
      TAKE addresses dataset distillation in the natural language domain, where the goal is to select a compact subset of a large text corpus that, when used for fine-tuning, recovers the performance of training on the full dataset.
      Influence-function-based selection methods — the dominant approach in this area — suffer from a hard-sample bias (HSB): samples that are difficult for the model to fit early in training receive inflated importance scores, leading to a selected corpus that over-represents noisy or boundary examples and suppresses the moderate examples that carry reliable, generalisable signal.
      TAKE corrects this bias by accumulating influence estimates across multiple checkpoints along the training trajectory rather than at a single snapshot, yielding trajectory-integrated knowledge scores whose bias properties are formally characterised in the full chapter.
      A synthetic candidate pool is then generated by fine-tuning a language model on the scored corpus, and a final compact subset is selected via discrete optimal transport, with a formal bound on the downstream loss gap introduced by this selection.
      TAKE outperforms DiLM and DaLLME on five benchmarks — AG News, IMDb, SST-2, MNLI-m, and QQP — and achieves a five-fold reduction in performance variance relative to a random selection baseline.
      This work has been submitted to ECML PKDD 2026.
    \end{foldable}

\end{enumerate}

\begin{foldable}
  Across all three contributions, a common set of principles is at work.
  Diversity is not treated as an implicit consequence of scale but as an explicit design target: each method introduces a dedicated mechanism — adaptive thresholds, contrastive priors, or trajectory-integrated scoring — specifically to measure and control the representational coverage of the data used for distillation.
  Practical constraints are taken as design inputs rather than obstacles: the black-box interface, the absence of real data, and the cost of large corpora are each incorporated directly into the method's formulation, ensuring that the solutions remain applicable in the deployment conditions that motivated them.
  Finally, each method pairs its data generation or selection step with a principled quality control procedure that filters or reweights the resulting pool, preventing low-quality synthetic or biased real examples from degrading the student's training signal.
  Together, these contributions demonstrate that diversity-aware distillation under restricted access is not only feasible but capable of matching or exceeding methods that operate under far less constrained conditions.
\end{foldable}
